Heterogeneous warehouse fleets combine agents with complementary roles, such as carrier vehicles and pickers, whose decisions must be coordinated under task dependencies, shared resources, and battery constraints. This thesis investigates whether locally deployed, off-the-shelf \glspl{llm} can support high-level decision-making for this form of multi-robot collaboration. A grounded control loop was implemented in the Battery-TA-RWARE simulation environment: the current warehouse state is translated into a prompt, the model selects discrete macro-actions, and proposed actions are validated against environment-defined feasibility constraints before execution.

The experiments evaluate six language models spanning approximately 1B to 14B parameters across six warehouse scenarios. They compare centralized joint planning with per-agent shared-context planning, natural-language with \gls{json}-based state representation, and three prompted objective families addressing shelf delivery, \gls{llm}-call usage, and battery management. Performance is assessed using shelf deliveries, model-call frequency, a call-normalised delivery measure, and charging behaviour. Generation and action-validation failures are logged for diagnostic run validation.

The results show that planning architecture, prompt configuration, model choice, and scenario configuration all influenced collaborative task performance. The natural-language configuration and shared-context planning were generally associated with stronger delivery performance in the evaluated setting, although the latter required more model interactions. Across the selected model artefacts, performance did not increase consistently with nominal parameter size, and changing the stated objective produced different behavioural responses across models.

These findings indicate that language models can produce grounded decisions within a constrained heterogeneous-robot control interface, while also showing that effective collaboration depends on the interaction between the model and the surrounding planning system. Given the limited number of runs and initialisations, the results are interpreted as descriptive evidence and motivate broader evaluation of language-model-based robot collaboration.
